For lighter atoms (and at higher levels), a substantial, or even dominant, part is played
\SourcePage{342}{345}
by nonradiative transitions. In such a transition, the energy released when an electron from a higher level fills the hole is used to eject another inner electron from the atom (the so-called \emph{Auger effect}); after this process the atom is left with two holes. To first order in $1/Z$, the probabilities of these processes and their contribution to the corresponding level width are independent of the atomic number (see the problem)\footnote{As an example, the Auger width of the $K$ level is about 1 eV, while for higher levels it reaches values of roughly 10 eV.}.
